\section{Proofs}\label{app:proofs}

\begin{proof}[Proof of Proposition~\ref{prop:cost}]
Take additive $U$. One item $a$ has score $1+\epsilon$ and cost $B$. There are $\lfloor B/c_{\min}\rfloor$ items of score $1$ and cost $c_{\min}$. The rule fills with $a$ first and then has no budget, so it earns $1+\epsilon$ against $\mathrm{OPT}\ge\lfloor B/c_{\min}\rfloor$. Let $\epsilon\to0$.
\end{proof}

\begin{proof}[Proof of Proposition~\ref{prop:recency}]
Walking from newest to oldest, keep each item that fits. Since the later items cost more than $B$ in total, one of them fails to fit, and at that moment the remaining budget is below its cost, which is at most $c(m)$. The remaining budget only shrinks afterwards, so $m$ never fits. A relevance rule takes $m$ first when it outranks those items, and $c(m)\le B$ then lets it in.
\end{proof}

\begin{proof}[Proof of Proposition~\ref{prop:stale}]
An invariant score cannot depend on how the two timestamps are assigned, so the ranking between the two versions is independent of which one is current. Under a uniform assignment the outdated version ranks first with probability $1/2$.
\end{proof}

\begin{proof}[Proof of Proposition~\ref{prop:selection}]
Sort the items by density $w/c$ and let $G$ be the prefix accepted before the first item $g$ that does not fit. If no item fails to fit, $G$ holds every item of positive weight and is optimal. Otherwise $G\cup\{g\}$ is a density-ordered prefix whose cost exceeds the budget, so its value is at least the optimum of the fractional relaxation, and hence at least $\mathrm{OPT}$. Thus $F(G)+w_g\ge\mathrm{OPT}$, and the larger of the two terms is at least $\frac12\mathrm{OPT}$. The density run contains $G$, so it has at least $F(G)$, and the best single item has at least $w_g$ because $c(g)\le B$.
\end{proof}

\begin{proof}[Proof of Proposition~\ref{prop:calibration}]
For any set $S$, $|\sum_{m\in S}p_m-\sum_{m\in S}\hat p_m|\le e(S)$. Hence
$\sum_{\hat S}p\ge\sum_{\hat S}\hat p-e(\hat S)\ge\alpha\sum_{S^\star}\hat p-e(\hat S)\ge\alpha\sum_{S^\star}p-\alpha e(S^\star)-e(\hat S)$,
where the middle step uses that $S^\star$ is feasible. For $\alpha=1$, the items in $S^\star\cap\hat S$ cancel: apply the same chain to $S^\star\setminus\hat S$ and $\hat S\setminus S^\star$, using $\sum_{\hat S\setminus S^\star}\hat p\ge\sum_{S^\star\setminus\hat S}\hat p$, which holds because $\hat S$ is optimal for $\hat p$.
\end{proof}

\begin{proof}[Proof of Proposition~\ref{prop:blind}]
Let each round have $n$ items with equal features, $\mu=1$, $c=1$, $B=1$ and $F(S;w)=\sum_{i\in S}w_i$, with the index of the valid item uniform and independent across rounds. What a blind learner sees in past rounds is independent of the current valid index, so whichever item it picks is valid with probability $1/n$ and $\mathbb{E}F_t(S_t)=1/n$, while $\mathrm{OPT}_t=1$. A learner that sees the version chain picks the current item in every round.
\end{proof}

\begin{proof}[Proof of Theorem~\ref{thm:rd}]
$\Phi$ is monotone submodular (facility location), $\Phi(\emptyset)=0$ and $\Phi(X)=1$. Let $g_i$ be the density-maximal pick, with gain $\Delta_i=\Phi(Z_i)-\Phi(Z_{i-1})$. Monotonicity and submodularity give $\Phi(Z^\star)\le\Phi(Z^\star\cup Z_{i-1})\le\Phi(Z_{i-1})+\sum_{z\in Z^\star\setminus Z_{i-1}}\Delta(z\mid Z_{i-1})\le\Phi(Z_{i-1})+R^\star\Delta_i/c(g_i)$. So $d_i:=\Phi(Z^\star)-\Phi(Z_i)\le(1-c(g_i)/R^\star)d_{i-1}\le e^{-c(g_i)/R^\star}d_{i-1}$ whenever $d_{i-1}>0$, using $1-x\le e^{-x}$ for all real $x$. If $d_{i-1}\le0$ for some $i$, then $\Phi(Z_\ell)\ge\Phi(Z^\star)$ by monotonicity. In every case $\Phi(Z_\ell)\ge(1-e^{-R_\ell/R^\star})\Phi(Z^\star)$, and substituting $\Phi=1-D$ gives the claim.
\end{proof}

\begin{proof}[Proof of Theorem~\ref{thm:store}]
Condition on $r$ and let $T$ be the time until the first gap of length $\tau$ in a Poisson process of rate $r$ started at the arrival. Let $A\sim\mathrm{Exp}(r)$ be the first access. If $A>\tau$ then $T=\tau$, otherwise $T=A+T'$ with $T'$ an independent copy of $T$ by the memoryless property. So $\mathbb{E}T=\mathbb{E}\min(A,\tau)+(1-e^{-r\tau})\mathbb{E}T$. With $\mathbb{E}\min(A,\tau)=(1-e^{-r\tau})/r$, this gives $\mathbb{E}T=(e^{r\tau}-1)/r$.

The pairs (arrival time, lifetime) form a marked Poisson process with i.i.d.\ marks. By the marking theorem, the number of items alive at $t$ is Poisson with mean $a\int_0^t\Pr[L>u]\,du$, which converges to $a\,\mathbb{E}L$ under the finiteness assumption.
\end{proof}

\section{Earlier tests}\label{app:past}
Three earlier evaluations of older releases shaped the method. Their result files remain in the repository, and their repositories later joined the fit.

\begin{center}\small
\resizebox{\textwidth}{!}{\begin{tabular}{@{}lll@{}}\toprule
Repositories & Finding & Effect on the method \\\midrule
5 Python & Near-twin rule took look-alike commit notes for versions of a fact & Only facts replace each other \\
8 Python & With keys ahead of every rival, without keys level with hybrid & Twins must cover the query \\
9 TypeScript, Go, Rust & No lead for a model fitted on Python alone (\V{old-kanerva_vs_hybrid-mean}, \V{olci-kanerva_vs_hybrid-mean}) & Those repositories joined the fit \\
\bottomrule\end{tabular}}
\end{center}
